\def\topic{Ordered pair}
\ifdefined\MyHeaderIncluded
\else
\def\qwertyflag{true}

\ifdefined\MyHeaderIncluded

\else
  \newcommand{\foobar}{FooBar}
  \def\MyHeaderIncluded{true}
% if is till the end of the file

\documentclass[10pt,a4paper]{article}
\usepackage[utf8]{inputenc}
%\usepackage[german]{babel}
%\usepackage[cyrillic]{babel}
\usepackage[T1]{fontenc}
%\usepackage[T2A]{fontenc}
%\usepackage[english, russian]{babel}
\usepackage{amsmath}
\usepackage{amsfonts}
\usepackage{amssymb}
\usepackage{amsthm}
\usepackage{mathtools} %coloneqq

\usepackage{bussproofs} %

%\renewcommand{\thesubsection}{\arabic{subsection}}
%\newtheorem{thm}{Theorem}[section]
%\renewcommand{\thethm}{\arabic{subsection}.\arabic{thm}}

\usepackage{thmtools}
\usepackage{hyperref}
\hypersetup{
    colorlinks=true,
    linkcolor=blue,
    filecolor=magenta,      
    urlcolor=cyan,
    pdftitle={Overleaf Example},
    %pdfpagemode=FullScreen,
    }
%\usepackage{stix-type1}%\turnediota
\theoremstyle{definition}

\newtheorem{constructor}{Constructor}[section]
\newtheorem{iaxiom}{Axiom}
\newtheorem{axiom}{Axiom}
\newtheorem*{axiom*}{Axiom}
\newtheorem{definition}{Definition}[section]
\newtheorem*{definition*}{Definition}
\newtheorem{remark}{Remark}[section]
\newtheorem*{remark*}{Remark}
\newtheorem{recall}{Recall}[section]
\newtheorem*{recall*}{Recall}

\newtheorem{lemma}{Lemma}[section]
\newtheorem{theorem}{Theorem}[section]
\newtheorem*{theorem*}{Theorem}
\newtheorem{note}{Note}[section]
\usepackage{etoolbox}

%\newcommand*{\@Dashv}[2]{%
%  \reflectbox{$\m@th#1#2$}%
%}
\newcommand*{\separation}[3]{% 1)binder, 2)set, 3)condition
  \{#1\in #2: #3\}
}

% https://www.freecodecamp.org/news/simple-regex-tricks-for-beginners-3acb3fa257cb/
%  \\mathcal\{T\} |---> \topology
%  \\mathcal\{([^\}]*)\}   |--->  $1
\newcommand{\AND}{\ \ \&\ \ }
\newcommand{\any}{\mathrm{\ any}}
\newcommand{\borel}{\mathcal{B}}
\newcommand{\can}{\mathrm{can}}
\newcommand{\canonize}{\mathrm{canonize}}
\newcommand{\closure}{\mbox{cl}}
\newcommand{\complex}{\mathbb{C}}
\newcommand{\converges}{\longrightarrow}
\newcommand{\countable}{\mbox{\ countable}}
\newcommand{\ldefined}{\mbox{\ defined}}
\newcommand{\dom}{\mbox{dom}}
\newcommand{\Dom}{\mbox{Dom}}
\newcommand{\Done}{\noindent \textbf{Done.}\\}
\newcommand{\lexists}{\mbox{\ exists}}
\newcommand{\lforall}{\mbox{\ forall}}
\newcommand{\lunique}{\mbox{\ unique}}
\newcommand{\family}{\mathrm{\ family}}
\newcommand{\Fm}{\mathrm{Fm}}
\newcommand{\FuncSymb}{\mathrm{FuncSymb}}
\newcommand{\function}{\mathrm{\ function}}
\newcommand{\functional}{\mbox{\ functional}}
\newcommand{\IFF}{\ \Longleftrightarrow\ }
\renewcommand{\iff}{\ \longleftrightarrow\ }
\newcommand{\IMPLIES}{\ \Longrightarrow\ }
\renewcommand{\implies}{\ \longrightarrow\ }
\newcommand{\Inf}{\mathrm{Inf}}
\newcommand{\injective}{\mbox{\ injective}}
\newcommand{\injection}{\mbox{\ injection}}
\newcommand{\interior}{\mbox{int}}
\newcommand{\intersection}{\bigcap}
\newcommand{\isempty}{\mathrm{\ empty}}
\newcommand{\LB}{\mathrm{LB}}
\newcommand{\nghbset}{\mathcal{N}}
\newcommand{\nonempty}{\mathrm{\ nonempty}}
\newcommand{\Ord}{\mathrm{Ord}}
\newcommand{\power}{\mathcal{P}}
\newcommand{\pr}{\mathrm{pr}}
\newcommand{\fst}{\mathrm{fst}}
\newcommand{\snd}{\mathrm{snd}}
\newcommand{\PredSymb}{\mathrm{PredSymb}}
\newcommand{\Proof}{\noindent \textbf{Proof.}\\}
\newcommand{\ran}{\mathrm{ran}}
\newcommand{\Ran}{\mathrm{Ran}}
\newcommand{\real}{\mathbb{R}}
\newcommand{\relation}{\mathrm{\ relation}}
\newcommand{\set}{\mathrm{\ set}}
\newcommand{\singleton}{\mathrm{\ singleton}}
\newcommand{\semantics}{\mathrm{semantics}}
\newcommand{\semialgebra}{\mathcal{S}}
\newcommand{\Sup}{\mathrm{Sup}}
\newcommand{\surjective}{\mbox{\ surjective\ on\ }}
\newcommand{\Term}{\mathrm{Term}}
\newcommand{\Type}{\mathrm{Type}}
\newcommand{\topology}{\mathcal{T}}
\newcommand{\total}{\mbox{\ total\ on\ }}
\newcommand{\transitive}{\mathrm{\ transitive}}
\newcommand{\union}{\bigcup}
\newcommand{\lTrue}{\mathrm{True}}
\newcommand{\lFalse}{\mathrm{False}}
\newcommand{\lUncertain}{\mathrm{Uncertain}}
\newcommand{\lif}{\mathrm{if}}
\newcommand{\lthen}{\mathrm{then}}
\newcommand{\lelse}{\mathrm{else}}
\newcommand{\lexception}{\mathrm{exception}}
\newcommand{\lundef}{\mathrm{undef}}
\newcommand{\univocal}{\mathrm{\ univocal}}
\newcommand{\univocalat}{\mathrm{\ univocal\ at\ }}
\newcommand{\UB}{\mathrm{UB}}
\newcommand{\Un}{\mathcal{U}}
\newcommand{\Var}{\mathrm{Var}}

\newcommand{\notvarphi}{\phi}
\renewcommand{\phi}{\varphi}

\newcommand{\comment}[1]{\quad\mbox{(#1)}}
% local commands
%\preto\section{\undef\blipo} %etoolbox
\preto\subsection{\undef\blipo} %etoolbox
% https://tex.stackexchange.com/questions/105684/limit-macro-newcommand-scope-to-section
\usepackage{graphicx}
\usepackage{authblk}

\setlength\parindent{0pt} % whole doc of \noindent
\usepackage{soul}
%\usepackage{mathgl}

\fi % guard
\ifdefined\topic
\else
\def\topic{No Titel}
\fi
\title{\topic}
\author[1]{Georgy Dunaev}
\begin{document}
\maketitle
\fi
\section{Ordered pairs of sets and classes}
\subsection{Introduction}
An ordered-pair construction must recover both components from the result.
For classes, there is an additional question: can either component be a
proper class? We first compare set-based constructions, then develop the
author's recursive tagged construction. The same notation $(a,b)$ is reused
locally; subscripts and subsection boundaries identify the construction in use.
The three displayed axioms are the intended interface, whose verification
is a separate task for each proposed model.
\subsection{Axioms of ordered pair}
\[
((a\in\Term)\land(b\in\Term))\implies((a,b)\in\Term)
\clabeltagd{axm}{DunaevTaggedOPairConstructor}
\]
\[
((a\set)\land(b\set))\implies((a,b)\set)
\clabeltagd{axm}{DunaevTaggedOPairSize}
\]%\label(axm:ordered pair set)
\[
((a_1,b_1)=(a_2,b_2))
\iff
(
(a_1=a_2)
\land
(b_1=b_2)
)
\clabeltagd{axm}{OPairUniqueness}
\]
\subsection{Kuratowski's ordered pair of sets}
For sets $a$ and $b$, the usual Kuratowski construction is
\begin{definition}[Kuratowski pair]
\[
(a,b)_K\ceq \{\{a,a\},\{a,b\}\}.
\]
\end{definition}
The inner collections distinguish the first component from the second.
The following comprehensions recover the singleton classes $\{a\}$ and
$\{b\}$; taking their unions recovers the components themselves.
\begin{definition}[Kuratowski projections]
\[
\fst_K(p)\ceq \bigcup\cc{x}{\forall y\in p,\ x\in y},
\qquad
\snd_K(p)\ceq \bigcup\cc{x}{\exists! y\in p,\ x\in y}.
\]
\end{definition}
These projection equations are intended for $p=(a,b)_K$ with set
components. The equal-component case is included: then $p=\{\{a\}\}$.

\begin{remark}[recall the definition]
\[
(\exists! y\in Q,\ \psi)
\stackrel{\mathrm{def}}{\IFF}
(\exists y\in Q,\ \psi)
\&
(! y\in Q,\ \psi)
\]
%\[
%\exists! y\in p ,\ x \in y
%\IFF
%(\exists y\in p,\ x \in y)
%\&
%(! y\in p,\ x \in y)
%\]
%because
%\[
%\exists (y\in p),\ (x \in y)
%\IFF
%\exists y,\ (y\in p) \land (x \in y)
%\]
%\[
%\! (y\in p),\ (x \in y)
%\IFF
%! y,\ (y\in p) \land (x \in y)
%\]
%\[
%\IFF
%\forall y_1, y_2,\ ((y_1\in p) \land (x \in y_1))\land
%(y_2\in p) \land (x \in y_2)
%\implies
%y_1 = y_2
%\]
%...
%\[
%\exists! y,\ (y\in p) \land (x \in y)
%\]
\end{remark}
\begin{theorem}[Kuratowski projections recover the components]
For sets $a,b$, $\fst_K((a,b)_K)=a$ and $\snd_K((a,b)_K)=b$.
\end{theorem}
\begin{proof}
If $a\ne b$, the elements common to $\{a\}$ and $\{a,b\}$ form
$\{a\}$, and the elements belonging to exactly one of these two sets form
$\{b\}$. If $a=b$, both comprehensions give $\{a\}$. Taking unions
proves the two identities in either case.
\end{proof}

\subsection{Definition of Metamath's ordered pair of sets}
The following convention deliberately restricts pairing to set components:
\begin{definition*}[Metamath's ordered pair] %maybe V somehow?
\[
\langle A,B \rangle=\cc{x}{(A\set)\land(B\set)\land (x\in\{\{A\},\{A,B\}\})}
\labeltagd{def}{MM_ordered pair}
\]
\end{definition*}
When a component is proper, the defining guard fails and the result is empty. Thus this convention serves set relations, but does not provide an injective pairing of arbitrary classes.
\subsection{Definition of Morse's ordered pair}
The next construction extends pairing to class arguments by first coding their elements. The products in this display use the preceding set-pair construction; the resulting class is then taken as a new pair:
\[
 s(x) = \{\emptyclass \} \cup \{\{t\} \mid t \in x\}
\]
\[
(x, y) = (\{0\} \times s(x)) \cup (\{1\} \times s(y))
\]
and also defines 3--tuple:
\[
(x, y, z) = (\{0\} \times s(x)) \cup (\{1\} \times s(y)) \cup (\{2\} \times s(z))
\]
Similar tuples which are presented in the section ``Classes as Types" are called lists in this text.\\
This approach distinguishes the auxiliary set pairs from the resulting class pairs, and correspondingly distinguishes their products and projections.\\

The following proposal instead seeks a uniform recursive definition for sets and classes.

\subsection{Recursive Dunaev tagged pairs}
\begin{remark}[comparison with the newer marker construction]
The Obelisk catalog's entry \texttt{DefDunaevOrderedPair} (3.76) uses a
separate marker-based encoding. Writing it as $D$ to avoid ambiguity,
\[
S_3(a)=\{\{\{a\}\}\},\qquad M(A,n)=\{\{0,n\}\}\cup\cc{\{S_3(a),n\}}{a\in A},
D(A,B)=M(A,1)\cup M(B,2).
\]
In contrast, this subsection develops the recursive pair $P$ given below.
Already $P(0,0)=\emptyclass$, whereas
$D(0,0)=\{\{0,1\},\{0,2\}\}$. The newer catalog's projection entries
3.77--3.78 concern $D$, so they should not be read as proofs of the
recursive construction. Both approaches are retained in their own setting.
\end{remark}

\paragraph{Local assumptions.}
We work with well-founded sets and use extensionality, separation,
replacement, and well-founded set recursion. These are available in the
usual ZF set part of NBG or MK. Class comprehension then extends the
construction to arbitrary classes. In particular, the argument below is not
an existence proof from extensionality and class comprehension alone.
No choice principle is needed.
The required membership-recursion theorem is proved in
\ref{thm:membership_recursion_on_W}; regularity identifies its domain
$\W$ with $\V$ in the present setting.

The simultaneous equations we want to justify are
\[
(a,b) \ceq (\{0\}\times a)\cup (\{1\}\times b),
\qquad
X\times Y \ceq \cc{(x,y)}{x\in X \land y\in Y}.
\]
The apparent mutual recursion disappears once the two tag maps are isolated.
We use $0=\emptyclass$ and $1=\{0\}$. Auxiliary graphs used to invoke the
recursion theorem may be coded by the already available Kuratowski pair;
this does not assume the new pair construction.

\begin{definition}[The two recursive tag maps]
On sets, define $R$ by well-founded recursion and then define $L$ by
\[
R(x)=\{0\}\cup\cc{R(u)}{u\in x},
\qquad
L(x)=\cc{R(u)}{u\in x}.
\label{def:recursive_pair_tags}
\]
Here each image is a set by replacement. The recursion step takes the
already constructed values on the members of $x$, forms their image, and
adjoins the empty set.
\end{definition}

\begin{lemma}[Existence and separation of the tags]
There is a unique class function $R:\V\to\V$ satisfying its displayed
recursion. For every set $x$, $0\in R(x)$ and $0\notin L(x)$. Consequently,
$L(x)\ne R(y)$ for all sets $x,y$.
\end{lemma}
\begin{proof}
Apply well-founded set recursion to the operation sending a set-sized
partial-value graph $g$ to $\{0\}\cup\operatorname{ran}(g)$. Its restriction
to the members of any set is a set graph by replacement, so the operation
has a set value at every stage. The recursion theorem gives existence;
membership induction gives uniqueness, since agreement on every member of
$x$ forces agreement at $x$.

The definition puts $0$ in every $R(u)$, so $R(u)$ is never empty.
Every member of $L(x)$ is some $R(u)$, hence is different from $0$.
Thus $L(x)$ omits the marker that every $R(y)$ contains.
\end{proof}

\begin{lemma}[Injectivity of each tag]
For sets $x,y$,
\[
R(x)=R(y)\iff x=y,\qquad L(x)=L(y)\iff x=y.
\label{lem:recursive_pair_tag_injectivity}
\]
\end{lemma}
\begin{proof}
For $R$, use membership induction on $x$ with the induction assertion
\[
\forall (y\set),\ R(x)=R(y)\implies x=y.
\]
Assume this assertion holds for every $u\in x$, and suppose $R(x)=R(y)$.
If $u\in x$, then $R(u)\in R(x)=R(y)$. Since $R(u)\ne0$, it is not the
adjoined marker; therefore $R(u)=R(v)$ for some $v\in y$. The induction
hypothesis for $u$, which quantifies over \emph{every} set $v$, gives $u=v$.
This proves $x\subseteq y$.

Conversely, if $v\in y$, then the nonempty set $R(v)$ belongs to $R(x)$,
so $R(v)=R(u)$ for some $u\in x$. The same induction hypothesis gives
$u=v$, proving $y\subseteq x$. Extensionality yields $x=y$ and completes
the induction. No induction hypothesis about the possibly unrelated set
$y$ was needed.

If $L(x)=L(y)$, adjoining $\{0\}$ to both sides gives $R(x)=R(y)$,
hence $x=y$. The reverse implications follow by substitution of equals.
\end{proof}

\begin{definition}[The recursive pair of classes]
For arbitrary classes $A,B$, put
\[
P(A,B)=\cc{L(a)}{a\in A}\cup\cc{R(b)}{b\in B}.
\label{def:recursive_class_pair}
\]
All members of $A$ and $B$ are sets, so the tags in this comprehension
have already been defined. No graph whose members have proper-class
components is being assumed. Write $(A,B)$ for $P(A,B)$ within this
construction, and use this pair in the associated Cartesian product.
\end{definition}

\begin{theorem}[The simultaneous equations have a unique solution]
The construction satisfies
\[
P(0,x)=L(x),\qquad P(1,x)=R(x)\quad(x\set),
\]
and, for arbitrary classes $A,B$,
\[
P(A,B)=\cc{P(0,a)}{a\in A}\cup\cc{P(1,b)}{b\in B}
       =(\{0\}\times A)\cup(\{1\}\times B).
\]
It is the unique class-valued operation satisfying these equations, with
set values on pairs of sets.
\end{theorem}
\begin{proof}
Since $L(0)=0$, the definition gives
$P(0,x)=\cc{R(u)}{u\in x}=L(x)$ and
$P(1,x)=\{L(0)\}\cup\cc{R(u)}{u\in x}=R(x)$.
Substituting these identities proves the equations for $P$ and the product.

For uniqueness, let $Q$ be another such operation. Its equations force
$Q(0,0)=0$ and then
\[
Q(1,x)=\{0\}\cup\cc{Q(1,u)}{u\in x}.
\]
Uniqueness of set recursion gives $Q(1,x)=R(x)$ for every set $x$.
The equation at $(0,x)$ now gives $Q(0,x)=L(x)$. Finally the equation
for arbitrary $Q(A,B)$ expresses it using these set tags, so $Q(A,B)=P(A,B)$.
\end{proof}

\begin{remark}
Do not confuse any of previous definitions with a family $F:2\to\mathcal{U}$.
\[
\tilde{K}\ceq \{\emptyclass\}\cup\cc{\{k\}}{k\in K}
\]
\[
[a,b] \ceq (\{0\}\times \tilde{a})\cup (\{1\}\times \tilde{b})
\]
Generally speaking,
\[
(a,b) :\neq [a,b]
\]
The extra empty marker in $\tilde K$ distinguishes an empty value from an absent fiber. The relation between this family encoding and the recursive pair construction requires its own verification.\\
\end{remark}

\subsection{Implementation sketch in Coq}
The following code illustrates the recursive construction for readers familiar with Coq. It is an excerpt, rather than a standalone proof certificate; its auxiliary definitions and proof obligations belong to the surrounding development.

The datatype $\texttt{Ens}$ represents sets in the ambient type theory. A construction on this datatype does not by itself establish the proper-class extension; that extension uses class comprehension in the mathematical development.

\begin{lstlisting}
(* The type of sets in a universe. *)
Inductive Ens : Type :=
    sup : forall A : Type, (A -> Ens) -> Ens.
\end{lstlisting}
\begin{lstlisting}
(* The empty set *)
Definition Vide : Ens :=
  sup False (fun x : False => match x return Ens with
                              end).
Definition Pair (A B:Ens) : Ens
 := sup bool (fun b : bool => if b then A else B).
(* The singleton set  *)
Definition Sing (E : Ens) := Pair E E.
\end{lstlisting}
\begin{lstlisting}
Definition Zero := Vide.
Definition One  := Sing Vide.
\end{lstlisting}
\begin{lstlisting}
Definition img (f:Ens -> Ens) (c:Ens) : Ens :=
match c with
| sup C h => sup _ (fun d => f (h d))
end.
\end{lstlisting}
\begin{lstlisting}
Definition bu A B := Union (Pair A B).
Infix "u" := bu (at level 60, right associativity).
\end{lstlisting}
The following straightforward implementation of the ordered pair is not accepted by the proof assistant, because the built--in termination checker does not accept the recursion as written.
\begin{lstlisting}
Fail
Program Fixpoint  OPair_Zero (c : Ens) : Ens
                    := bu (img OPair_Zero Zero) (img OPair_One c)
            with  OPair_One  (c : Ens) : Ens
                    := bu (img OPair_Zero One ) (img OPair_One c)
.
\end{lstlisting}
The first values help expose the recursion:
\begin{lstlisting}
Definition img_OPair_Zero_on_Zero := Vide.
(* (img OPair_Zero Zero) === {(0,k)|k in 0} = emptyset *)
Definition        OPair_Zero_Zero := Vide.
(* OPair Zero Zero === (0,0) = {(0,j)|j in 0}u{(1,k)|k in 0} = emptyset *)
Definition img_OPair_Zero_on_One  := Sing Vide.
(*(img OPair_Zero One) === {(0,k)|k in 1} = {(0,0)} = {emptyset}*)
Definition        OPair_Zero_One  := Sing (Sing Vide).
(* OPair Zero One === (0,1) = {} u {(1,0)} = {{emptyset}} *)


\end{lstlisting}
The mathematical reduction above uses only one recursive function. A
structural version of its right tag can be written as follows:
\begin{lstlisting}
Fixpoint OPair_One (c : Ens) {struct c} : Ens :=
  match c with
  | sup C h =>
      sup (unit + C)%type
        (fun i => match i with
                  | inl _ => Vide
                  | inr u => OPair_One (h u)
                  end)
  end.

Definition OPair_Zero (c : Ens) : Ens := img OPair_One c.
\end{lstlisting}
The left summand supplies the empty marker; the right summand recurses
on members of the input. This makes the decrease visible in the constructor
instead of hiding it inside a call to \texttt{img}. The excerpt illustrates
the construction; a mechanized certificate must additionally verify the
universe declarations and compatibility with the representation's
extensional equality. No compilation result for this excerpt is asserted
here.

The pair operation then uses the two images:
\begin{lstlisting}
Definition OPair (a b: Ens) : Ens := (img OPair_Zero a) u (img OPair_One b).
\end{lstlisting}

\subsection{Projections and the ordered-pair property}
The preceding construction explains both the recursion and the tags.
The projections now recover arbitrary class components by comprehension.

\begin{definition}[Recursive tagged projections]
For any class $C$, define
\[
\pr_1(C)=\cc{a}{L(a)\in C},\qquad
\pr_2(C)=\cc{b}{R(b)\in C}.
\label{def:recursive_pair_projections}
\]
These are the original proposed formulas
\[
\pr_1(C)=\cc{a}{(0,a)\in C},\qquad
\pr_2(C)=\cc{b}{(1,b)\in C},
\]
since comprehension ranges over set witnesses.
\end{definition}

\begin{theorem}[Projection identities]
For all classes $A,B$,
\[
\pr_1(P(A,B))=A,\qquad\pr_2(P(A,B))=B.
\label{thm:recursive_pair_projections}
\]
\end{theorem}
\begin{proof}
For a set $a$, membership $L(a)\in P(A,B)$ says either
$L(a)=L(u)$ for some $u\in A$, or $L(a)=R(v)$ for some $v\in B$.
The second alternative is excluded by the marker distinction. The first
is equivalent to $a\in A$ by injectivity of $L$. Thus the first projection
has exactly the members of $A$.

Likewise, $R(b)$ can belong only to the right tagged image, and injectivity
of $R$ makes $R(b)\in P(A,B)$ equivalent to $b\in B$. Class
extensionality gives both identities.
\end{proof}

\begin{theorem}[The ordered-pair property]
For arbitrary classes $A_1,B_1,A_2,B_2$,
\[
P(A_1,B_1)=P(A_2,B_2)
\iff (A_1=A_2)\land(B_1=B_2).
\label{thm:recursive_pair_uniqueness}
\]
\end{theorem}
\begin{proof}
Apply the two projections to the asserted pair equality. The projection
identities give equality of the first and of the second components.
The converse follows directly from the definition of $P$.
\end{proof}

\begin{theorem}[Sethood of a recursive pair]
For arbitrary classes $A,B$,
\[
(P(A,B)\set)\iff (A\set)\land(B\set).
\label{thm:recursive_pair_sethood}
\]
\end{theorem}
\begin{proof}
If $A,B$ are sets, replacement makes their two tagged images sets, and
pairing and union make their union a set.

Conversely, suppose $P(A,B)$ is a set. Separation gives the subset
\[
\cc{c\in P(A,B)}{\exists(a\set),\ c=L(a)}.
\]
Each of its members has a unique inverse image under $L$ by tag
injectivity. Replacement applied to this inverse relation therefore makes
$A$ a set. The same argument with $R$ makes $B$ a set. These inverse
images are uniquely determined, so choice is not needed.
\end{proof}

\paragraph{Why the height argument is unnecessary here.}
The recursive definition uses well-founded set recursion, and its injectivity
uses membership induction. Both are explicit assumptions of this local
construction. There is no need to assign a new numerical height to the pair
codes. In particular, the bounded depth measure discussed in the height
section need not be substituted for ordinal rank: its failure to decrease
strictly at limit values has no bearing on the tag-injectivity proof above.
In a setting without regularity, the same argument applies to its
well-founded part when the required recursion and replacement principles
are available; it makes no assertion about arbitrary non-well-founded sets.

\subsection{Finite examples and Cartesian products}
The recursive tags have particularly small initial values:
\[
L(0)=0,\quad R(0)=1,\quad L(1)=\{1\},\quad R(1)=2,
\]
where $0=\emptyclass$, $1=\{0\}$, and $2=\{0,1\}$ are von Neumann
ordinals. Thus the pair table is
\[
\begin{array}{c|ccc}
P & 0 & 1 & 2\\\hline
0 & 0 & \{1\} & \{1,2\}\\
1 & 1 & 2 & 3\\
2 & \{0,\{1\}\} & \{0,1,\{1\}\} & \{0,1,\{1\},2\}
\end{array}
\]
The row is the first component. For example,
\[
P(2,2)=\{L(0),L(1),R(0),R(1)\}
      =\{0,\{1\},1,2\}=\power 2.
\]
The number of members of a pair of finite sets is the sum of the numbers
of members of its components, because the two tag maps are injective and
have disjoint images. This does not make the pair operation ordinal
addition: $P(1,2)=3$, whereas $P(2,1)=\{0,1,\{1\}\}$ is not an ordinal.

The Cartesian product instead collects a pair for each choice of a member
of each factor. Its corresponding table is
\[
\begin{array}{c|ccc}
\times & 0 & 1 & 2\\\hline
0 & 0 & 0 & 0\\
1 & 0 & 1 & \{0,\{1\}\}\\
2 & 0 & 2 & \{0,1,\{1\},2\}
\end{array}
\]
Thus $P(2,2)=2\times2=\power2$ happens for this small example; pairing
and Cartesian product remain different operations.

On the truth-value representatives $\Prop=\{0,1\}$, the product table
restricts to Boolean conjunction: a product is $1$ exactly when both
factors are $1$. In particular,
\[
1\times1=\{P(0,0)\}=\{0\}=1.
\]
Intersection gives the same conjunction on these two representatives.

\ifdefined\qwertyflag
\undef\qwertyflag
\end{document}
\fi
